% Pile cuivre-argent et mesure de sa tension à vide (exercice 3, chapitre 8)
\begin{tikzpicture}
\pcPile{pcCuivre}{pcBleu}{pcArgent}{pcEau}
\draw[pcFil] (pileG) -- (-0.4,4.4) -- (6.6,4.4) -- (pileD);
\pcAppareilRond{(3.1,4.4)}{V}
\node[pcLegende,anchor=south east] at (2.72,4.45) {COM};
\node[pcLegende,anchor=south west] at (3.48,4.45) {V};
\node[pcLegendeG,anchor=south] at (3.1,4.85) {tension affichée : $+\qty{0,43}{V}$};
\node[pcLegende] at (3.1,2.3) {pont salin};
\node[pcLegende,anchor=north,align=center] at (0,-0.15) {lame de cuivre\\solution de nitrate de cuivre(II)};
\node[pcLegende,anchor=north,align=center] at (6.2,-0.15) {lame d'argent\\solution de nitrate d'argent};
\end{tikzpicture}
\par\smallskip\legende{Le voltmètre mesure la tension à vide de la pile ; sa borne COM est reliée à la lame de cuivre.}
